\documentclass[a4paper,12pt]{report}

% ============================================================
%  PACKAGES
% ============================================================
\usepackage[utf8]{inputenc}
\usepackage[T1]{fontenc}
\usepackage[english]{babel}
\usepackage[margin=2.5cm]{geometry}
\usepackage{graphicx}
\usepackage{xcolor}
\usepackage{colortbl}
\usepackage{longtable}
\usepackage{array}
\usepackage{booktabs}
\usepackage{titlesec}
\usepackage{fancyhdr}
\usepackage{tikz}
\usepackage{listings}
\usepackage{hyperref}
\usepackage{enumitem}
\usepackage{tocloft}
\usepackage{caption}
\usepackage{float}

\definecolor{primaryblue}{RGB}{16,42,87}
\definecolor{accentgreen}{RGB}{124,206,6}
\definecolor{lightgray}{RGB}{245,245,245}

\titleformat{\chapter}[display]
  {\normalfont\huge\bfseries\color{primaryblue}}
  {\chaptertitlename\ \thechapter}{20pt}{\Huge}
\titleformat{\section}{\large\bfseries\color{primaryblue}}{\thesection}{1em}{}
\titleformat{\subsection}{\normalsize\bfseries}{\thesubsection}{1em}{}

\pagestyle{fancy}
\fancyhf{}
\fancyhead[L]{\small\textit{Final Year Project Report}}
\fancyhead[R]{\small\thepage}
\renewcommand{\headrulewidth}{0.4pt}

\newcolumntype{P}[1]{>{\raggedright\arraybackslash}p{#1}}

\lstset{
  basicstyle=\ttfamily\small,
  breaklines=true,
  backgroundcolor=\color{lightgray},
  frame=single,
  numbers=left,
  numberstyle=\tiny,
  keywordstyle=\color{primaryblue}\bfseries,
  commentstyle=\color{gray}\itshape,
}

% ============================================================
\begin{document}

% ---------------- TITLE PAGE ----------------
\begin{titlepage}
\vspace*{-1cm}

%-------------------------------------------------
% Top section: Left logo - Center text - Right logo
%-------------------------------------------------
\noindent
\begin{minipage}[t]{0.2\textwidth}
    \raggedright
    \includegraphics[height=3cm]{images/tekup.png}
\end{minipage}
\hfill
\begin{minipage}[t]{0.56\textwidth}
    \centering
    \textbf{Republic of Tunisia}\\
    \textbf{Ministry of Higher Education and Scientific Research}\\[0.4cm]
    \textbf{\Large Private Higher School of Engineering and Technology}
\end{minipage}
\hfill
\begin{minipage}[t]{0.2\textwidth}
    \raggedleft
    \includegraphics[height=3cm]{images/BYB.png}
\end{minipage}

\vspace{1cm}

\begin{center}

{\Large\textbf{Final Year Project Report}}\\[0.4cm]

Submitted in partial fulfillment of the requirements for the\\
\textbf{National Engineering Degree in Applied Sciences and Technology}\\[0.3cm]

Major: Software Engineering and Information Systems\\[1cm]

Prepared by\\[0.2cm]
{\Large\textbf{Med Aziz Selmi}}\\[1.2cm]

%=================================================
% Project title inside decorative lines
%=================================================

{\color{primaryblue}\rule{\textwidth}{6pt}}\\
{\color{primaryblue}\rule{\textwidth}{2pt}}

\vspace{0.6cm}

{\LARGE\bfseries Treyo ---}\\[0.25cm]
{\LARGE\bfseries Design and Development of an Intelligent}\\[0.25cm]
{\LARGE\bfseries Trainer--Learner Matching Platform}

\vspace{0.6cm}

{\color{primaryblue}\rule{\textwidth}{2pt}}\\
{\color{primaryblue}\rule{\textwidth}{6pt}}

\vspace{1.2cm}

\end{center}

\begin{flushleft}
Client organization: \hfill LeanConsulting\\[0.3cm]
Industry supervisor: \hfill Mrs. Nour Jlassi, Designer\\[0.3cm]
Academic supervisor: \hfill Mr. Med Amine Ben Rhouma, Assistant Professor
\end{flushleft}

\vfill

\begin{center}
{\large\textbf{Academic Year 2025--2026}}
\end{center}

\end{titlepage}

\pagenumbering{roman}

% ---------------- DEDICATION ----------------
\chapter*{Dedication}
\addcontentsline{toc}{chapter}{Dedication}
\emph{To be filled in as you wish.}

% ---------------- ACKNOWLEDGEMENTS ----------------
\chapter*{Acknowledgements}
\addcontentsline{toc}{chapter}{Acknowledgements}
I would like to express my sincere gratitude to the team at
\textbf{LeanConsulting} for their trust and for providing the resources and
business context necessary for the success of this project. I would also
like to thank my industry supervisor and my academic supervisor for their
guidance, valuable advice, and availability throughout this work, as well
as the members of the jury for the interest they bring to it.

% ---------------- TABLE OF CONTENTS ----------------
\tableofcontents
\listoffigures
\listoftables

% ---------------- GENERAL INTRODUCTION ----------------
\chapter*{General Introduction}
\addcontentsline{toc}{chapter}{General Introduction}
\pagenumbering{arabic}

On-demand learning continues to grow, driven by the digitalization of
professional training. However, most existing platforms follow a top-down
model in which educational content is published directly by the publisher
or administrator, with no genuine quality validation of content proposed by
trainers, and no mechanism for forming learner groups based on actual
demand.

The \textbf{Treyo (BYB\_Match)} project addresses this gap by proposing a
platform that connects \textbf{trainers} and \textbf{learners}, organized
around a two-level administrative validation flow: trainers are approved
by an administrator upon registration, and every course they submit ---
attached to a module created by the administrator --- is then reviewed,
priced, and published (or rejected) by that same administrator before it
becomes visible to learners. Training groups are formed automatically as
soon as enough learners express interest in a course, and a machine
learning--based recommendation engine personalizes content discovery.

This report presents the design and development approach of the platform,
organized according to the \textbf{Scrum} agile methodology. The first
chapter sets out the general context of the project and the study of
existing solutions. The second chapter covers requirements analysis and
specification as well as the overall architecture. The following chapters
detail each \emph{release} and the sprints that make it up. The report
concludes with a general conclusion and future perspectives.

% ============================================================
\chapter{Project Framework}
% ============================================================

\section{General project framework}

\subsection{Presentation of the host organization}
This project was carried out on behalf of \textbf{LeanConsulting}, which
expressed the need for a digital platform to structure and digitalize its
trainer--learner matching activity. LeanConsulting acts here as the end
client of the product: the requirements specification, the validation of
deliverables at the end of each sprint, and functional acceptance testing
were carried out in coordination with its representatives, while the
development itself was conducted as part of the Final Year Project at
[school name].

\subsection{Context presentation}
This project was carried out as part of the Final Year Project in Software
Engineering and Information Systems, in response to the need expressed by
LeanConsulting. It consists of designing and developing, end to end, an
on-demand, multi-actor training platform, marketed under the name
\textbf{Treyo}.

\subsection{Context and objectives of the project}
The main objective is to provide an environment where:
\begin{itemize}[leftmargin=*]
  \item the \textbf{administrator} structures the catalog by creating
        \emph{modules} (categories) and by approving trainer registrations
        as well as the courses they submit;
  \item the \textbf{trainer}, once approved, proposes their own courses
        under an existing module, attaches training material, and awaits
        validation and pricing by the administrator;
  \item the \textbf{learner} discovers published courses --- via search or
        personalized recommendation --- expresses interest, joins an
        automatically formed group, pays online, follows scheduled
        sessions, communicates through messaging, and leaves a review at
        the end of the course.
\end{itemize}

Secondary objectives include: setting up a multi-channel notification
system, a comprehensive admin dashboard, a recommendation service based on
user interactions, and transparent financial management (daily trainer
compensation, online learner payment).

\section{Study of existing solutions}

\subsection{Study of the Udemy platform}
Udemy allows any trainer to publish a course in a self-service manner,
without prior qualitative validation of the educational content by a third
party. The model relies on pre-recorded, freely accessible videos, with no
cohort formation or real-time group interaction.

\subsection{Study of the Preply platform}
Preply connects learners and trainers for private (1-to-1) lessons, with
integrated payment and messaging. However, the platform does not offer
automatic group formation based on a demand threshold, nor an
administrative content-validation flow managed by a central curator.

\subsection{Study of the Superprof platform}
Superprof operates as a matching directory with no content-validation
workflow and no integrated centralized payment; negotiation and payment
take place outside the platform.

\section{Comparison of the solutions studied}

\begin{table}[H]
\centering
\caption{Comparison of existing platforms}
\begin{tabular}{|P{3cm}|c|c|c|c|}
\hline
\rowcolor{primaryblue!15}
\textbf{Criterion} & \textbf{Udemy} & \textbf{Preply} & \textbf{Superprof} & \textbf{Treyo} \\
\hline
Content validation by an administrator & No & No & No & \textbf{Yes} \\
\hline
Automatic group formation & No & No & No & \textbf{Yes} \\
\hline
Integrated payment & Yes & Yes & No & \textbf{Yes} \\
\hline
Personalized recommendation (ML) & Yes & Partial & No & \textbf{Yes} \\
\hline
Transparent daily trainer compensation & No & No & No & \textbf{Yes} \\
\hline
Group messaging & No & Yes (1-to-1) & No & \textbf{Yes} \\
\hline
\end{tabular}
\end{table}

\section{Synthesis for the choice of solution}
None of the solutions studied combine an administrative content-validation
flow for trainers, demand-driven group formation, and fine-grained
financial traceability (per-course pricing, transparent daily trainer
compensation). This combination constitutes Treyo's differentiating
positioning and justifies custom development rather than adopting an
existing solution.

\section{Project management methodology adopted}

\subsection{Scrum methodology}
The project was managed according to the \textbf{Scrum} agile methodology.
The product backlog was split into \textbf{4 releases}, themselves broken
down into one- to two-week \textbf{sprints}, each delivering a tested,
demonstrable functional increment.

\subsection{Roles in Scrum}
\begin{itemize}[leftmargin=*]
  \item \textbf{Product Owner}: owns the business need, prioritizes the
        backlog.
  \item \textbf{Scrum Master}: facilitates the process, removes obstacles.
  \item \textbf{Development team}: designs, codes, and tests every
        increment.
\end{itemize}

\subsection{Design methodology adopted}
The design follows the \textbf{UML} formalism (use case, class, and
sequence diagrams) for each sprint, complemented by the technical
layered architecture described in the following chapter.

% ============================================================
\chapter{Requirements Analysis and Specification}
% ============================================================

\section{Requirements specification}

\subsection{Functional requirements}
\begin{itemize}[leftmargin=*]
  \item \textbf{Account management}: registration, login, email
        verification, password reset, trainer approval by the
        administrator.
  \item \textbf{Module management}: creation, editing, and archiving of
        course categories by the administrator.
  \item \textbf{Course management}: creation by the trainer (title,
        description, duration, level, format, prerequisites, PDF/PPT
        material), validation/rejection and pricing by the administrator.
  \item \textbf{Group formation and enrollment}: expressing interest,
        automatic group formation at the required threshold, session
        scheduling.
  \item \textbf{Payment}: online payment via the Konnect (ClickToPay)
        gateway, confirmation via a verified webhook.
  \item \textbf{Messaging}: group conversation between learners and the
        trainer, with the administrator as an implicit member of every
        group.
  \item \textbf{Reviews}: learner ratings and reviews at the end of a
        course, moderated by the administrator.
  \item \textbf{Notifications}: push and in-app notifications on key
        events (approval, group formation, message).
  \item \textbf{Recommendation}: personalized course suggestions via a
        machine learning service.
  \item \textbf{Trainer compensation}: daily rate set by the administrator
        per course, viewed by the trainer through a monthly earnings
        dashboard.
  \item \textbf{Admin dashboard}: platform-wide statistics, management of
        users, courses, groups, reviews, and system settings (revenue
        display currency).
\end{itemize}

\subsection{Non-functional requirements}
\begin{itemize}[leftmargin=*]
  \item \textbf{Security}: JWT-based authentication, password hashing,
        role-based authorization (\texttt{STUDENT}, \texttt{TRAINER},
        \texttt{ADMIN}), account-enumeration prevention on sensitive
        endpoints.
  \item \textbf{Performance}: API response time under 300\,ms on common
        read endpoints.
  \item \textbf{Availability}: decoupled (stateless API) architecture
        allowing horizontal scalability.
  \item \textbf{Internationalization}: mobile interface available in
        French, English, Arabic (with RTL support), and Spanish.
  \item \textbf{Portability}: cross-platform mobile application (iOS /
        Android) via Expo/React Native.
  \item \textbf{Maintainability}: layered architecture, clear separation
        between controllers, services, and the persistence layer.
\end{itemize}

\section{Global use case diagram}
Figure~\ref{fig:uc-global} presents the three main actors (\emph{Student},
\emph{Trainer}, \emph{Administrator}) and their high-level use cases.

\begin{figure}[H]
\centering
% \includegraphics[width=0.9\textwidth]{uc_global.png}
\caption{Global use case diagram}
\label{fig:uc-global}
\end{figure}

\section{Global class diagram}
The domain model is built on an abstract \texttt{BaseEntity} class
(identifier, creation and update timestamps) inherited by all persistent
entities. The main entities are: \texttt{User} (abstract), specialized into
\texttt{Student}, \texttt{Trainer}, \texttt{Admin}; \texttt{CourseModule};
\texttt{Course}; \texttt{Group}; \texttt{Enrollment}; \texttt{Payment};
\texttt{CourseReview}; \texttt{Conversation} and \texttt{Message};
\texttt{Notification}; \texttt{Interaction}; \texttt{SearchLog};
\texttt{SystemSetting}. Statuses (\texttt{ApprovalStatus},
\texttt{GroupStatus}, \texttt{EnrollmentStatus}, \texttt{PaymentStatus},
etc.) are modeled as typed enumerations rather than free-form strings,
ensuring consistent state transitions.

\begin{figure}[H]
\centering
% \includegraphics[width=0.95\textwidth]{class_diagram_global.png}
\caption{Global class diagram}
\label{fig:class-global}
\end{figure}

\section{Work planning}

\subsection{Release breakdown}
\begin{table}[H]
\centering
\caption{Release breakdown}
\begin{tabular}{|c|P{7cm}|c|}
\hline
\rowcolor{primaryblue!15}
\textbf{Release} & \textbf{Content} & \textbf{Sprints} \\
\hline
Release 1 & Authentication, account management, and onboarding & Sprint 1, 2 \\
\hline
Release 2 & Modules, course creation and validation & Sprint 3, 4 \\
\hline
Release 3 & Group formation, scheduling, payment & Sprint 5, 6 \\
\hline
Release 4 & Messaging, reviews, notifications, ML recommendation, admin dashboard, trainer compensation & Sprint 7, 8 \\
\hline
\end{tabular}
\end{table}

\subsection{Product backlog}
The complete backlog comprises 60 user stories spread across the 4
releases above. Table~\ref{tab:backlog-extrait} shows a representative
excerpt; the full backlog is provided in the appendix.

\begin{table}[H]
\centering
\caption{Excerpt from the product backlog}
\label{tab:backlog-extrait}
\begin{tabular}{|c|P{9cm}|c|}
\hline
\rowcolor{primaryblue!15}
\textbf{ID} & \textbf{User Story} & \textbf{Priority} \\
\hline
US-13 & As a trainer, I want to create a course under an existing module so that I can propose training content. & High \\
\hline
US-17 & As an administrator, I want to set the course price and currency at approval time so that pricing policy is centralized. & High \\
\hline
US-26 & As the system, I want to automatically form a group once the enrollment threshold is reached so that training can start. & High \\
\hline
US-32 & As a learner, I want to pay for a course via Konnect so that I can confirm my enrollment. & High \\
\hline
US-49 & As a trainer, I want to see my current month's earnings on my home screen so that I have a quick financial overview. & High \\
\hline
\end{tabular}
\end{table}

\section{Application architecture}

\subsection{Logical architecture}
The application follows a layered architecture spread across four
independent components communicating via REST APIs:

\begin{itemize}[leftmargin=*]
  \item \textbf{Backend} (Spring Boot 3.2, Java 17): exposes the main REST
        API, organized in \texttt{Controller} $\rightarrow$
        \texttt{Service} $\rightarrow$ \texttt{Repository} layers (Spring
        Data JPA / Hibernate) on top of a \textbf{PostgreSQL} database.
        Security is handled by Spring Security and JWT tokens (the
        \texttt{jjwt} library).
  \item \textbf{Mobile application} (React Native / Expo SDK 54,
        TypeScript): used by learners and trainers, organized by route
        groups (\texttt{expo-router}) according to the logged-in role.
  \item \textbf{Admin dashboard} (Next.js 16, React 19, TypeScript,
        Tailwind): web interface reserved for the administrator.
  \item \textbf{Recommendation service} (Python, FastAPI): independent
        microservice leveraging interaction and search logs to compute
        personalized course recommendations.
\end{itemize}

\begin{figure}[H]
\centering
% \includegraphics[width=0.85\textwidth]{architecture_logique.png}
\caption{Logical architecture of the platform}
\label{fig:archi-logique}
\end{figure}

\subsection{Physical architecture}
The four components are deployed independently and communicate
exclusively via HTTP/JSON: the mobile application and the admin dashboard
consume the Spring Boot API, which in turn queries PostgreSQL and calls
the recommendation microservice for personalized discovery endpoints. The
Konnect payment gateway communicates with the backend through a
server-to-server webhook, verified on the backend side before any
enrollment is confirmed.

\begin{figure}[H]
\centering
% \includegraphics[width=0.85\textwidth]{architecture_physique.png}
\caption{Physical deployment architecture}
\label{fig:archi-physique}
\end{figure}

\section{Project mockups}
Low-fidelity then high-fidelity mockups were produced for the key user
journeys (registration, login, course creation, dashboard) before
implementation, in order to validate the user experience with
stakeholders.

\begin{figure}[H]
\centering
% \includegraphics[width=0.3\textwidth]{maquette_register.png}
\caption{Mockup --- registration screen}
\end{figure}

\section{Difficulties encountered}
\begin{itemize}[leftmargin=*]
  \item Naming conflict between the business entity \texttt{Module} and
        the JDK's \texttt{java.lang.Module} class, causing Lombok
        compilation errors on unrelated classes; resolved by renaming the
        entity to \texttt{CourseModule}.
  \item Coordinating the file-upload flow before the \texttt{Course}
        entity exists (the training material must be uploaded before the
        course is created); resolved via a dedicated endpoint keyed on the
        trainer's identifier rather than the course's.
  \item Preventing account enumeration on the forgot-password and
        resend-verification endpoints, requiring a strictly identical
        response whether the email exists or not.
  \item Computing the trainer's daily earnings from a semi-structured JSON
        session schedule, with deduplication of days containing multiple
        sessions of the same course.
\end{itemize}

\section{Working environment}

\subsection{Project management tools}
Backlog and sprint tracking via a shared Kanban board; version control
with \textbf{Git}, hosted on \textbf{GitHub}.

\subsection{UML modeling tools}
\textbf{StarUML} for class, use case, and sequence diagrams.

\subsection{Development and testing framework}
\begin{table}[H]
\centering
\caption{Technologies used per component}
\begin{tabular}{|P{3.5cm}|P{9cm}|}
\hline
\rowcolor{primaryblue!15}
\textbf{Component} & \textbf{Technologies} \\
\hline
Backend & Java 17, Spring Boot 3.2.2, Spring Security, Spring Data JPA / Hibernate, JWT (jjwt 0.12.3), Maven \\
\hline
Database & PostgreSQL \\
\hline
Mobile application & React Native 0.81, Expo SDK 54, TypeScript, expo-router, i18next, Axios \\
\hline
Admin dashboard & Next.js 16.2, React 19, TypeScript, Tailwind CSS, Recharts \\
\hline
Recommendation service & Python, FastAPI \\
\hline
Payment & Konnect (ClickToPay) \\
\hline
Push notifications & Expo Notifications \\
\hline
Transactional email & SMTP (Gmail) \\
\hline
\end{tabular}
\end{table}

\subsection{Development environment}
IDE: IntelliJ IDEA (backend), VS Code / WebStorm (mobile and web). API
tooling: Postman. API documentation: Swagger / OpenAPI (springdoc).

\subsection{Operating systems}
Development was carried out on Windows 11, with cross-testing on Android
and iOS emulators via Expo Go.

% ============================================================
\chapter{Release 1: Authentication and Account Management}
% ============================================================

\section{Sprint 1: Registration, login, and roles}

\subsection{Sprint 1 objectives}
Set up the authentication foundation (registration, login, email
verification, password reset) and the distinction between the three
application roles.

\subsection{Sprint 1 backlog}
\begin{itemize}[leftmargin=*]
  \item US-01 --- Student registration.
  \item US-03 --- Token-based email verification.
  \item US-04 --- Password reset (without account enumeration).
  \item US-05 --- Login and JWT token issuance.
  \item US-06 --- Authenticated password change.
\end{itemize}

\subsection{Functional requirements specification}
An unauthenticated user can create a student or trainer account. A trainer
account remains in \texttt{PENDING} status until administrative validation
(handled in Sprint 2). Any failed login attempt returns a generic error
message, without indicating whether the email exists.

\subsection{Class diagram}
\begin{figure}[H]
\centering
% \includegraphics[width=0.7\textwidth]{class_sprint1.png}
\caption{Class diagram --- Authentication}
\end{figure}

\subsection{Dynamic diagrams}
\begin{figure}[H]
\centering
% \includegraphics[width=0.7\textwidth]{seq_login.png}
\caption{Sequence diagram --- Login}
\end{figure}

\subsection{Implementation}
Screenshots of the registration, login, and profile management screens.

\subsection{Technologies used}
Spring Security, JWT, BCrypt, SecureStore (mobile) for token storage.

\section{Sprint 2: Onboarding and trainer approval}

\subsection{Sprint 2 objectives}
Implement the multi-step onboarding flow (student and trainer) and the
administrative trainer-approval workflow, with email notification in both
cases (approval / rejection).

\subsection{Sprint 2 backlog}
\begin{itemize}[leftmargin=*]
  \item US-09 --- Student onboarding (interests, education level).
  \item US-10 --- Trainer onboarding (specializations, CV).
  \item US-07 --- Admin review of trainer applications.
  \item US-08 --- Approval / rejection email.
\end{itemize}

\subsection{Functional requirements specification}
The trainer uploads a CV and fills in their specializations, skills, and
experience. The administrator reviews this information in the
\emph{Pending Approvals} queue and makes a decision, systematically
triggering an asynchronous email (best-effort, never blocking the API
response on a delivery failure).

\subsection{Class diagram}
\begin{figure}[H]
\centering
% \includegraphics[width=0.7\textwidth]{class_sprint2.png}
\caption{Class diagram --- Onboarding and approval}
\end{figure}

\subsection{Implementation}
Onboarding screens (mobile) and trainer-application review interface
(admin dashboard).

% ============================================================
\chapter{Release 2: Modules and Course Lifecycle}
% ============================================================

\section{Sprint 3: Module management}

\subsection{Sprint 3 objectives}
Allow the administrator to structure the catalog into modules
(categories) that trainers use to classify their courses.

\subsection{Sprint 3 backlog}
\begin{itemize}[leftmargin=*]
  \item US-12 --- Module CRUD (creation, editing, archiving).
\end{itemize}

\subsection{Functional requirements specification}
A module has a unique name and a description. Archived modules disappear
from the trainer-facing picker without affecting courses already attached
to them.

\subsection{Class diagram}
\begin{figure}[H]
\centering
% \includegraphics[width=0.6\textwidth]{class_sprint3.png}
\caption{Class diagram --- CourseModule}
\end{figure}

\subsection{Implementation}
\emph{Modules} page of the admin dashboard: listing, creation, editing,
archive/unarchive.

\section{Sprint 4: Course creation and validation}

\subsection{Sprint 4 objectives}
Implement course creation by the trainer (with an educational attachment)
and the validation/rejection/pricing workflow by the administrator,
replacing the previous model in which the administrator directly created
and assigned courses to trainers.

\subsection{Sprint 4 backlog}
\begin{itemize}[leftmargin=*]
  \item US-13 --- Course creation by the trainer.
  \item US-14 --- Material upload (PDF/PPT/ZIP) before the course exists.
  \item US-15 --- Editing a course in \texttt{PENDING}/\texttt{REJECTED} status.
  \item US-16 --- Admin review queue.
  \item US-17 --- Setting the price and currency at approval time.
  \item US-18 --- Rejection with an optional note.
  \item US-20 --- Filtering out non-approved courses on the public side and in ML.
\end{itemize}

\subsection{Functional requirements specification}
The trainer submits a course with \texttt{PENDING} status and no price
field: pricing is an exclusively administrative decision, made at approval
time. Rejection is possible with a note sent to the trainer by email, who
can then fix the course and resubmit it.

\subsection{Class diagram}
\begin{figure}[H]
\centering
% \includegraphics[width=0.75\textwidth]{class_sprint4.png}
\caption{Class diagram --- Course and ApprovalStatus}
\end{figure}

\subsection{Dynamic diagrams}
\begin{figure}[H]
\centering
% \includegraphics[width=0.7\textwidth]{seq_approve_course.png}
\caption{Sequence diagram --- Course approval by the administrator}
\end{figure}

\subsection{Implementation}
Mobile course-creation screen (trainer) with module picker and material
upload; \emph{Pending Courses} dashboard page (review, pricing,
approval/rejection).

% ============================================================
\chapter{Release 3: Group Formation and Payment}
% ============================================================

\section{Sprint 5: Group formation and scheduling}

\subsection{Sprint 5 objectives}
Automate group formation once the required enrollment threshold is
reached, and allow the trainer to schedule sessions.

\subsection{Sprint 5 backlog}
\begin{itemize}[leftmargin=*]
  \item US-25 --- Expressing interest in a course.
  \item US-26 --- Automatic group formation.
  \item US-27 --- Session scheduling by the trainer.
  \item US-28 --- Schedule viewing on the learner side.
  \item US-31 --- Automatic closure of completed groups/enrollments.
\end{itemize}

\subsection{Functional requirements specification}
Each expression of interest is counted against the course's
\texttt{minStudentsRequired} threshold. A scheduled job (plus a catch-up
pass on application startup) closes groups whose end date has passed.

\subsection{Class diagram}
\begin{figure}[H]
\centering
% \includegraphics[width=0.75\textwidth]{class_sprint5.png}
\caption{Class diagram --- Group and Enrollment}
\end{figure}

\subsection{Dynamic diagrams}
\begin{figure}[H]
\centering
% \includegraphics[width=0.7\textwidth]{seq_group_formation.png}
\caption{Sequence diagram --- Automatic group formation}
\end{figure}

\subsection{Implementation}
Join-request button, session-scheduling screen (trainer), \emph{My
Schedule} screen (learner).

\section{Sprint 6: Online payment}

\subsection{Sprint 6 objectives}
Integrate the Konnect (ClickToPay) payment gateway to confirm paid
enrollments.

\subsection{Sprint 6 backlog}
\begin{itemize}[leftmargin=*]
  \item US-32 --- Initiating a payment.
  \item US-33 --- Confirmation webhook verification.
  \item US-34 --- Payment history on the learner side.
\end{itemize}

\subsection{Functional requirements specification}
Although unauthenticated (server-to-server), the Konnect webhook is never
trusted blindly: every incoming event is revalidated against the Konnect
API before any enrollment is activated.

\subsection{Dynamic diagrams}
\begin{figure}[H]
\centering
% \includegraphics[width=0.7\textwidth]{seq_payment.png}
\caption{Sequence diagram --- Payment and webhook confirmation}
\end{figure}

\subsection{Implementation}
Mobile payment flow (Konnect redirect), payment history page.

% ============================================================
\chapter{Release 4: Messaging, Reviews, Notifications, and Oversight}
% ============================================================

\section{Sprint 7: Messaging, reviews, and recommendation}

\subsection{Sprint 7 objectives}
Add group messaging, the course review system, and the personalized
recommendation engine.

\subsection{Sprint 7 backlog}
\begin{itemize}[leftmargin=*]
  \item US-35 --- Group conversation.
  \item US-38 --- Course review at the end of the course.
  \item US-22 --- Personalized recommendations (ML service).
\end{itemize}

\subsection{Functional requirements specification}
Every group has an associated conversation, of which the administrator is
an implicit member. Every review updates the rolling average of the
course and the trainer. The recommendation service systematically excludes
non-approved courses from its results; its internal design and algorithms
are detailed in full in Chapter~\ref{chap:ml}.

\subsection{Class diagram}
\begin{figure}[H]
\centering
% \includegraphics[width=0.8\textwidth]{class_sprint7.png}
\caption{Class diagram --- Conversation, Message, CourseReview}
\end{figure}

\subsection{Implementation}
Group messaging screen, review form, \emph{Recommended for you} section.

\section{Sprint 8: Notifications, admin dashboard, and trainer compensation}

\subsection{Sprint 8 objectives}
Finalize multi-channel notifications, enrich the admin dashboard
(statistics, currency settings), and set up daily trainer-compensation
tracking.

\subsection{Sprint 8 backlog}
\begin{itemize}[leftmargin=*]
  \item US-40 --- Push notifications on key events.
  \item US-45 --- Revenue display currency (admin setting).
  \item US-48 --- Daily trainer rate per course.
  \item US-49 --- Current month earnings on the trainer home screen.
  \item US-50 --- Day-by-day earnings breakdown, browsable by month.
  \item US-51 --- Global statistics on the admin dashboard.
\end{itemize}

\subsection{Functional requirements specification}
The administrator sets, for each course, a daily amount compensating the
trainer, independent of the price paid by the learner. The earnings
computation service iterates over the trainer's groups, parses the JSON
session schedule, and aggregates one amount per unique day and per course
for the requested month; the associated mobile screen resets automatically
every new month since it always queries the current year and month.

\subsection{Class diagram}
\begin{figure}[H]
\centering
% \includegraphics[width=0.75\textwidth]{class_sprint8.png}
\caption{Class diagram --- Notification, SystemSetting, Course.trainerDailyRevenue}
\end{figure}

\subsection{Dynamic diagrams}
\begin{figure}[H]
\centering
% \includegraphics[width=0.7\textwidth]{seq_earnings.png}
\caption{Sequence diagram --- Computing the trainer's monthly earnings}
\end{figure}

\subsection{Implementation}
Admin dashboard (statistics, user and course management, \emph{Trainer
revenue / day} column), \emph{Settings} page (currency), mobile
\emph{Earnings} screen (trainer) with a month picker and a detailed
day-by-day list.

% ============================================================
\chapter{The Recommendation Engine (Machine Learning)}
\label{chap:ml}
% ============================================================

The personalized course recommendation feature is handled by a dedicated
microservice, entirely decoupled from the Spring Boot backend, so that
model training and inference never compete with the transactional API for
resources. This chapter details its design, algorithms, and operational
concerns, which the earlier sprint-level description in Sprint~7
intentionally kept at a high level.

\section{Goals and design constraints}
\begin{itemize}[leftmargin=*]
  \item \textbf{Cold start}: a brand-new learner has no interaction
        history, so the engine must still return relevant results from
        their onboarding profile alone.
  \item \textbf{Domain correctness over raw score}: a learner registered
        in \emph{Computer Science} must never see a highly-rated
        \emph{Marketing} course ranked above a mediocre in-domain one —
        domain match is treated as a hard constraint, not just a signal.
  \item \textbf{Freshness without downtime}: newly approved courses and
        newly logged interactions must eventually influence
        recommendations without ever taking the service offline to
        retrain.
  \item \textbf{Low latency under load}: recommendation calls are
        interactive (mobile home screen, search results), so the engine
        must answer in milliseconds even as the course catalog grows.
\end{itemize}

\section{Technology stack}
\begin{table}[H]
\centering
\caption{Recommendation engine technology stack}
\begin{tabular}{|P{3.5cm}|P{9cm}|}
\hline
\rowcolor{primaryblue!15}
\textbf{Layer} & \textbf{Technology} \\
\hline
API layer & Python, FastAPI, Uvicorn \\
\hline
Data access & psycopg2 (\texttt{ThreadedConnectionPool}) against the same PostgreSQL instance as the backend \\
\hline
Data manipulation & pandas \\
\hline
Content-based filtering & scikit-learn \texttt{TfidfVectorizer}, \texttt{cosine\_similarity} \\
\hline
Approximate nearest-neighbor search & FAISS (\texttt{IndexFlatIP}), with an automatic sklearn fallback if FAISS is unavailable \\
\hline
Collaborative filtering & scikit-learn \texttt{cosine\_similarity} over a user--item interaction matrix \\
\hline
\end{tabular}
\end{table}

\section{Data sources}
At startup, and on every scheduled refresh, the engine loads five
dataframes directly from PostgreSQL:
\begin{itemize}[leftmargin=*]
  \item \texttt{students} --- active students with their declared
        \texttt{primary\_domains}, \texttt{specific\_interests}, and
        \texttt{experience\_level};
  \item \texttt{courses} --- restricted to courses that are active,
        published, \textbf{approved by the administrator}
        (\texttt{approval\_status IS NULL OR approval\_status = 'APPROVED'}),
        whose trainer is active, and whose trainer has spare capacity
        (\texttt{active\_group\_count < max\_concurrent\_groups}), computed via
        a \texttt{WITH} CTE joining the \texttt{groups} table;
  \item \texttt{interactions} --- every logged student--course interaction
        (view, favorite, request, enroll, etc.);
  \item \texttt{trainers} --- active trainers with specializations,
        skills, and rating;
  \item \texttt{search\_logs} --- the last 10{,}000 search queries, used
        as an additional personalization signal.
\end{itemize}

This is the exact mechanism by which the admin approval workflow
(Chapter~\ref{chap:ml}'s upstream dependency, Release~2) reaches the
recommendation engine: a course sitting in \texttt{PENDING} or
\texttt{REJECTED} status is excluded at the SQL level and can never be
recommended, regardless of how well it matches a learner's profile.

\section{Content-based filtering}
Each course is reduced to a single \texttt{content\_text} field
concatenating its title, description, domain, and topic, then vectorized
with a TF-IDF vectorizer (\texttt{max\_features=1000},
English stop-words removed). Rows are L2-normalized so that inner product
equals cosine similarity, and indexed in a FAISS \texttt{IndexFlatIP}
structure for exact nearest-neighbor search in $O(\log N)$ rather than the
$O(N)$ cost of a brute-force sklearn comparison — the previous
implementation additionally used to materialize an $N \times N$
course-to-course similarity matrix, which was dropped as unnecessary
memory overhead. A free-text query (a learner's interests, or a search
string) is embedded the same way and scored against every course.

\section{Collaborative filtering}
A user--item interaction matrix is built from the \texttt{interactions}
table. For a given learner, the engine finds the most behaviourally
similar students via cosine similarity over that matrix, then surfaces
courses those neighbours engaged with that the target learner has not yet
seen.

\section{Hybrid ranking}
The production recommendation function (\texttt{hybrid\_recommendations})
blends four signals with fixed weights:

\begin{table}[H]
\centering
\caption{Hybrid recommendation source weights}
\begin{tabular}{|P{5cm}|c|P{6cm}|}
\hline
\rowcolor{primaryblue!15}
\textbf{Source} & \textbf{Weight} & \textbf{Signal} \\
\hline
Content-based & 45\% & Domain + interest match to course content \\
\hline
Collaborative & 35\% & Interactions of behaviourally similar students \\
\hline
Search history & 15\% & The learner's own past search queries \\
\hline
Popularity & 5\% & Highly-rated trainers within the learner's declared domain \\
\hline
\end{tabular}
\end{table}

Scores from all four sources are aggregated per course
(\texttt{groupby(course\_id)}, weighted sum), then re-sorted with a
\textbf{two-level sort key}: \texttt{in\_domain} (descending) first, raw
aggregated score (descending) second. This guarantees that no
out-of-domain course can ever outrank an in-domain one, even if its raw
collaborative or search score is numerically higher — domain is treated
as an explicit, authoritative signal, while collaborative/search are only
inferred ones.

\section{Cold-start strategy}
The engine selects its strategy per learner based on how many logged
interactions they have:

\begin{table}[H]
\centering
\caption{Recommendation strategy by interaction count}
\begin{tabular}{|c|P{9cm}|}
\hline
\rowcolor{primaryblue!15}
\textbf{Interaction count} & \textbf{Strategy} \\
\hline
0 & Cold start: rank courses purely by the learner's declared interests, domains, and experience level (from onboarding). \\
\hline
1--5 & Content-based filtering only. \\
\hline
6+ & Full hybrid (content + collaborative + search + popularity). \\
\hline
\end{tabular}
\end{table}

For interaction-count 0, if the learner's row is missing from the
in-memory snapshot (e.g. they just registered after the last model
refresh), the engine transparently falls back to a live, single-row
database query rather than returning an empty result.

\section{Explainability}
Every recommendation carries a short human-readable reason string (e.g.
\emph{"Matches your interests • Highly rated (4.6/5.0)"}), generated by
inspecting which signals contributed to that specific course's score. A
separate \texttt{explain\_recommendation} endpoint returns a fuller
breakdown (confidence score, contributing factors) for a given
learner/course pair, primarily used for debugging and product analysis
rather than end-user display.

\section{Model freshness and concurrency}
Rather than reloading the model on every request or blocking behind a
lock during retraining, the engine uses an \textbf{atomic swap} pattern:

\begin{itemize}[leftmargin=*]
  \item A background job periodically calls \texttt{rebuild\_models()},
        which re-fetches all five dataframes from PostgreSQL and builds a
        brand-new \texttt{CourseRecommendationSystem} instance
        \emph{without touching} the live one.
  \item Once the new model is fully built, all six state fields (data
        frames + model + version counter) are swapped in under a single
        lock acquisition (\texttt{\_swap\_in}), so in-flight requests that
        already took a snapshot of the old model finish normally, and any
        new request sees the new model immediately.
  \item If the rebuild fails for any reason, the exception is caught and
        the previous model stays active — a broken refresh never takes
        the service down.
\end{itemize}

Every request first calls \texttt{\_snapshot()} to atomically read
\texttt{(model, students, courses, interactions, trainers, searches,
version)} under a short-lived lock, then operates purely on those local
references, immune to a concurrent background swap.

\section{Caching layer}
A thread-safe, bounded LRU cache sits in front of the recommendation
computation, keyed by \texttt{(student\_id, n\_recommendations)}. A cache
entry is invalidated when any of three conditions is met:
\begin{enumerate}[leftmargin=*]
  \item the model version has advanced past the entry's version (a
        background refresh happened);
  \item the learner's interaction count has changed since the entry was
        cached (they clicked/enrolled/favorited something new);
  \item a defensive TTL (15 minutes by default) has elapsed.
\end{enumerate}
Whenever a new interaction is tracked via \texttt{track\_interaction()},
the learner's cache entries are explicitly invalidated so their very next
call reflects the fresh signal, even before the next scheduled model
refresh. Cache size and TTL are both configurable via environment
variables so capacity can be tuned in production without a code change.

\section{Trending and popularity endpoints}
Beyond personalized recommendations, the service also exposes a trending-
courses endpoint, ranking either by domain-scoped popularity or, in the
absence of a domain filter, by a simple \texttt{(num\_enrolled, rating)}
descending sort over the same approved-course dataset.

\section{Summary}
The recommendation engine is the component of the platform most directly
responsible for content \emph{discovery} quality: it is the layer that
converts the administrator's approval decisions (Release~2) and the
learners' behavioural signals (Interaction, SearchLog) into a ranked,
explainable, and continuously refreshed list of course suggestions, while
guaranteeing zero-downtime retraining and strict domain relevance.

% ============================================================
\chapter{Security}
\label{chap:security}
% ============================================================

Security concerns were addressed throughout every release rather than as
a single retrofit; this chapter consolidates them into one place.

\section{Authentication: JWT}
Authentication is stateless, backed by signed JSON Web Tokens
(HS256, via the \texttt{jjwt} library). On a successful login or
registration, \texttt{JwtService} issues:
\begin{itemize}[leftmargin=*]
  \item a short-lived \textbf{access token} embedding the user's
        \texttt{userId} and \texttt{role} as custom claims, signed with
        an HMAC key derived from a server-side secret (\texttt{jwt.secret},
        externalized to a gitignored \texttt{application.properties} file,
        never committed);
  \item a longer-lived \textbf{refresh token} carrying only the subject
        (email), used to renew the access token.
\end{itemize}
Every incoming request passes through a single
\texttt{JwtAuthenticationFilter}, registered \emph{before}
\texttt{UsernamePasswordAuthenticationFilter} in the Spring Security
filter chain, which extracts and validates the bearer token, then
populates the security context with the resolved role. Session creation
policy is set to \texttt{STATELESS} --- the server keeps no session state
at all, which is what allows the API to scale horizontally without a
shared session store.

\section{Authorization: role-based access control}
Authorization is enforced centrally in \texttt{SecurityConfig}, not
scattered across individual controllers, via a URL-pattern matcher:

\begin{table}[H]
\centering
\caption{Access control rules}
\begin{tabular}{|P{6cm}|P{5cm}|}
\hline
\rowcolor{primaryblue!15}
\textbf{Route pattern} & \textbf{Access} \\
\hline
\texttt{/api/auth/**}, \texttt{/api/files/download/**}, \texttt{/api/trainers}, \texttt{/api/modules}, \texttt{/api/settings/revenue-currency} & Public (no token required) \\
\hline
\texttt{/api/payments/konnect-webhook} & Public, but re-verified server-to-server (see below) \\
\hline
\texttt{/api/admin/**} & \texttt{ADMIN} role only \\
\hline
Everything else & Any authenticated user (role checked in-controller where finer granularity is needed, e.g. a trainer editing only their own course) \\
\hline
\end{tabular}
\end{table}

\section{Password storage}
Passwords are never stored or logged in plaintext. \texttt{BCryptPasswordEncoder}
(Spring Security's default, adaptive-cost hashing) is used both to hash a
new password at registration and to verify a login attempt, so the
plaintext password never touches persistent storage at any point in its
lifecycle.

\section{Account-enumeration prevention}
The \texttt{/api/auth/forgot-password} and
\texttt{/api/auth/resend-verification} endpoints are a common attack
surface for harvesting valid registered emails (a different response for
"email exists" vs. "email doesn't exist" leaks that information). Both
endpoints in Treyo return an \textbf{identical HTTP 200 response
regardless of whether the email is registered}, and the mobile UI shows
the same "if an account exists, check your inbox" message either way ---
so an attacker gains nothing by probing the endpoint with a list of
guessed addresses.

\section{Payment webhook integrity}
The Konnect (ClickToPay) webhook endpoint is necessarily unauthenticated
--- it is called server-to-server by Konnect's infrastructure, which
cannot present a Treyo-issued JWT. To prevent a forged webhook call from
fabricating a paid enrollment, the handler never trusts the incoming
payload directly: on every event it makes a synchronous call back to
Konnect's own payment-status API to confirm the transaction actually
succeeded before activating anything. An attacker who discovers the
webhook URL and POSTs a fake "payment succeeded" body gains nothing, since
the backend independently re-checks the payment's real status.

\section{File-upload handling}
Trainer CVs and course materials (PDF/PPT/ZIP) are restricted by MIME type
at the mobile document picker and re-validated server-side before being
written to disk under a namespaced path (keyed by trainer ID rather than a
client-supplied file name), preventing path traversal and reducing the
risk of arbitrary file execution on the server.

\section{Secrets management}
All credentials --- database password, JWT signing secret, Gmail SMTP
password, Konnect API keys, Gemini API keys --- live in gitignored files
(\texttt{application.properties}, \texttt{.env}) with only \texttt{.example}
templates tracked in version control, so a cloned or forked repository
never leaks production secrets.

\section{Cross-Origin Resource Sharing (CORS)}
CORS is centrally configured to allow all origins by pattern during
development (the admin dashboard and mobile app are served from different
hosts/ports), restricted to \texttt{GET/POST/PUT/DELETE/PATCH/OPTIONS},
with the \texttt{Authorization} header explicitly exposed so the frontend
can read it after a cross-origin response.

% ============================================================
\chapter{Notifications and the Social Feed}
\label{chap:notif-feed}
% ============================================================

\section{Push notification delivery}
Rather than integrating separately with Firebase Cloud Messaging (Android)
and Apple Push Notification service (iOS), the backend sends every push
through \textbf{Expo's push notification HTTP endpoint}
(\texttt{https://exp.host/--/api/v2/push/send}), which forwards to the
correct native platform on Treyo's behalf. This removes the need to
provision and rotate native push certificates on the backend server
itself.

\subsection{Device registration}
Each user device registers an Expo push token, persisted in a
\texttt{DeviceToken} entity keyed by a unique token value and indexed by
\texttt{userId}. A user can hold multiple tokens (e.g. phone \emph{and}
tablet); conversely, if the same physical device later logs into a
different account, the registration call overwrites the token's owner so
notifications never leak to the previous user.

\subsection{Sending}
\texttt{PushNotificationService} exposes \texttt{sendToUser} and the bulk
variant \texttt{sendToUsers} (used for group-wide events such as a new
chat message). Every send is annotated \texttt{@Async} so the triggering
request (e.g. posting a chat message) never blocks on the round-trip to
Expo's servers, and failures are logged rather than propagated --- a push
delivery failure must never fail the underlying business operation.

\section{In-app notification center}
Independently of push delivery, every notification is also persisted as a
\texttt{Notification} row so the user can review their notification
history inside the app, with read/unread state, even if the push itself
was missed (app closed, permissions denied, etc.).

\section{The AI-generated social feed}
The home-screen feed shown to students is not user-generated content ---
it is a daily batch of short posts generated by \textbf{Google Gemini}
through \texttt{FeedGenerationService}, designed to keep the app engaging
between training sessions.

\subsection{Generation pipeline}
\begin{itemize}[leftmargin=*]
  \item A scheduled job (\texttt{@Scheduled}, firing daily at 06:00
        server time) sends one structured prompt to Gemini asking for
        exactly 4 posts, mixing categories: one \texttt{FUN\_FACT}, two
        \texttt{DOMAIN\_TIP} posts (each covering a different training
        domain), and one \texttt{MOTIVATION} post.
  \item The prompt constrains title length (3--8 words), body length
        (30--60 words), tone (conversational, no clichés), and forbids
        emojis, and requires strict JSON output so the response can be
        parsed directly without scraping markdown code fences.
  \item All users see the \emph{same} daily batch --- a deliberate choice
        that keeps moderation simple (one thing to review, not N
        personalized streams) and creates a shared community feed rather
        than isolated per-user content.
  \item If Gemini returns malformed JSON or the API quota is exhausted,
        the service falls back to a small hard-coded seed pool, so the
        feed is never empty in production.
  \item A manual \texttt{generateBatch} endpoint exists for testing and to
        backfill the feed immediately after a fresh deployment, without
        waiting for the next scheduled run.
\end{itemize}

\subsection{Translation}
After the English batch is generated, the same service requests
translations into French, Arabic, and Spanish (stored as
\texttt{FeedPostTranslation} rows keyed by post and language code), so the
feed respects the learner's chosen app language without a second round
of AI generation per language selection.

\subsection{Moderation and reactions}
Users can react to feed posts (like/engagement toggle). Because content
is centrally generated rather than user-submitted, there is no
user-generated-content moderation queue for the feed --- moderation
effort is instead spent reviewing the daily Gemini prompt and output
format.

% ============================================================
\chapter{The Admin Dashboard}
\label{chap:admin-dashboard}
% ============================================================

The admin dashboard is a separate Next.js application, reserved
exclusively for the \texttt{ADMIN} role, and is the nerve center of the
platform's day-to-day operation. Beyond the Modules and Pending Courses
pages already detailed in Release~2, the dashboard exposes the following
pages.

\section{Dashboard overview}
The landing page (\texttt{/dashboard}) surfaces platform-wide KPIs as stat
cards (total users, students, trainers, enrollments, revenue) plus three
charts built with Recharts: a user-distribution pie chart
(students vs. trainers), a course-status pie chart (live vs. deleted), an
enrollments bar chart (active vs. completed), and a "this week's activity"
bar chart (new users, new enrollments, new courses). A quick-stats footer
row surfaces secondary counters (active users today, total interactions,
total messages, total groups).

\section{User management}
Lists every student and trainer, with search/filter and a detail view.
Includes the \emph{Pending Approvals} queue described in Release~1
(Sprint~2), where the administrator reviews a trainer's onboarding
submission (specializations, CV) before approving or rejecting the
account.

\section{Course management}
The \emph{Courses} page (Release~2, extended in Release~4) lists every
live course with: level badge, enrolled-student count, an inline-editable
minimum-students-required field, price, and --- added in the final
iteration of the project --- a \emph{Trainer revenue / day} column letting
the administrator set or edit the daily compensation rate directly from
the table, without navigating to a separate screen. The former
"Requested" column (interest-funnel count) was removed from this table
once a dedicated \emph{Requests} page (below) took over that
responsibility, to avoid duplicating the same metric in two places.

\section{Requests}
Shows the interest funnel per course (how many learners requested it vs.
how many are actually enrolled), letting the administrator gauge demand
before a group is even formed, and includes a cleanup action for stale,
never-converted join requests.

\section{Group management}
Lists every group across the platform with its status (forming, ready,
active, completed, cancelled), current/maximum size, and assigned
trainer, giving the administrator visibility into cohort formation
health without needing to inspect individual courses.

\section{Reviews moderation}
Every learner review (rating + written comment) is listed here. An
administrator can toggle a review's visibility to hide inappropriate
content from public display; the underlying rating value still counts
toward the course's and trainer's rolling average even when hidden, so
moderation cannot be used to inflate ratings by selectively hiding
negative-but-legitimate reviews.

\section{Notifications composer}
Lets the administrator send platform-wide or audience-targeted
notifications (e.g. "all students in the Web Development module"),
reusing the same \texttt{PushNotificationService} / \texttt{Notification}
persistence path described in Chapter~\ref{chap:notif-feed}.

\section{System health}
A live status page surfacing basic operational signals (API availability,
recent error indicators) so the administrator has a first line of
diagnosis before escalating to the development team.

\section{Settings}
Introduced in the platform's final iteration: a single \emph{Revenue
currency} setting, backed by a generic key--value \texttt{SystemSetting}
entity, letting the administrator choose the currency (TND, USD, EUR,
GBP, MAD, DZD, EGP, SAR, AED) in which the dashboard's aggregate revenue
figure is displayed. The setting is read publicly
(\texttt{GET /api/settings/revenue-currency}) so any screen --- including
the trainer-facing pricing form during course approval --- can default to
the same currency, while only the administrator can change it.

% ============================================================
\chapter{Internationalization}
\label{chap:i18n}
% ============================================================

\section{Supported languages}
The mobile application is fully localized in four languages --- French,
English, Arabic, and Spanish --- via \texttt{i18next} /
\texttt{react-i18next}, with per-language JSON dictionaries under
\texttt{i18n/locales/}. The active language is persisted in
\texttt{AsyncStorage} so it survives an app restart, and can be changed at
any time from \emph{Settings~$\rightarrow$~Language}.

\section{Right-to-left (RTL) support}
Arabic requires the interface to mirror horizontally (right-to-left
reading and layout direction). React Native's \texttt{I18nManager} handles
most of this automatically by flipping \texttt{flexDirection} and
absolute-positioned \texttt{left}/\texttt{right} styles across the entire
component tree once RTL mode is enabled --- switching to or from Arabic
therefore requires an app reload (via \texttt{expo-updates}) for the new
layout direction to fully apply, and the language-settings screen prompts
the user accordingly.

\subsection{A layout edge case: the background gradients}
One subtlety surfaced during Arabic testing: the app's ambient background
(four \texttt{LinearGradient} layers glowing from each screen edge) is
positioned with absolute \texttt{left}/\texttt{right} offsets. Under RTL,
React Native auto-mirrors those offsets --- moving each glow to the
opposite physical edge --- but does \emph{not} mirror the gradient's own
internal color direction, producing a visible seam where the two effects
disagreed. The fix wraps the entire background layer in a container with
an explicit \texttt{direction: 'ltr'} style, which opts that subtree out
of RN's automatic RTL mirroring while leaving the actual screen content
(text, buttons, navigation) in normal RTL flow. This keeps the background
visually identical across all four languages while ensuring reading order,
navigation icons, and text alignment still correctly mirror for Arabic
users.

% ============================================================
\chapter{Testing, Quality Assurance, and Deployment}
\label{chap:testing}
% ============================================================

\section{Testing strategy}
Given the project's four-component architecture, testing effort was
distributed according to where defects were most likely and most costly:
business-logic-heavy backend services and the recommendation engine's
external contract.

\subsection{Backend}
The Spring Boot application includes a context-load smoke test
(\texttt{BackendApplicationTests}), verifying the full Spring context ---
every controller, service, and repository bean, plus the JPA/Hibernate
mapping layer --- wires up correctly against the configured database on
every build. Business logic was additionally verified through manual,
scenario-driven testing at the end of every sprint (see each sprint's
\emph{Realisation} section), using Postman collections against the
Swagger/OpenAPI-documented endpoints.

\subsection{Recommendation microservice}
The ML service ships its own black-box HTTP test script
(\texttt{test\_api.py}), exercising the running FastAPI service exactly as
a real client would, end to end:
\begin{itemize}[leftmargin=*]
  \item \texttt{test\_health} --- confirms the model is loaded and reports
        student/course/interaction counts;
  \item \texttt{test\_recommendations} --- requests personalized
        recommendations for a known student and checks the response
        shape;
  \item \texttt{test\_cold\_start} --- verifies the interest-based ranking
        path used for zero-interaction learners;
  \item \texttt{test\_track\_interaction} --- posts a new interaction and
        confirms it is accepted;
  \item \texttt{test\_trending} --- checks the popularity/trending
        endpoint;
  \item \texttt{test\_explain} --- verifies the explanation endpoint
        returns a human-readable justification for a given
        recommendation.
\end{itemize}
This test script doubles as executable documentation of the service's
public contract and was run after every change to the ranking algorithm
or its weighting scheme, to catch regressions before they reached the
mobile app's recommendation feed.

\subsection{Mobile and admin dashboard}
The mobile application and the admin dashboard were validated through
manual functional testing on real devices/emulators (iOS and Android via
Expo Go) and browsers, following the acceptance criteria defined in each
user story of the product backlog, plus continuous TypeScript strict-mode
type-checking (\texttt{tsc --noEmit}) on every change as a first line of
defense against integration errors between the frontend and the evolving
backend API contracts.

\section{Deployment topology}
At the current stage of the project, the four components run as
independent local/staging processes: the Spring Boot backend on a JVM
process backed by PostgreSQL, the FastAPI recommendation service as a
separate Uvicorn process against the same database, the Next.js admin
dashboard served via its own Node process, and the mobile application
distributed for testing through Expo Go (LAN/tunnel modes) rather than a
signed store build.

\section{Path to production}
The following steps remain to move from the current development setup to
a production deployment, and are called out explicitly as they fall
outside the scope of what was completed during the six months of this
internship:
\begin{itemize}[leftmargin=*]
  \item Containerizing the backend, ML service, and admin dashboard
        (Docker) and defining a reverse-proxy / orchestration layer.
  \item Setting up a continuous-integration pipeline to run the backend
        smoke test and the ML service's black-box test suite on every
        push.
  \item Producing a signed production mobile build via \textbf{EAS
        Build}, with FCM/APNs credentials registered in the Expo project
        so push notifications work outside of Expo Go.
  \item Externalizing configuration (database URL, JWT secret, Konnect
        and Gemini API keys, SMTP credentials) through environment
        variables or a secrets manager rather than local property files.
\end{itemize}

% ============================================================
\chapter*{General Conclusion}
\addcontentsline{toc}{chapter}{General Conclusion}
% ============================================================

This project made it possible to design and develop a complete
trainer--learner matching platform, structured around a two-level
administrative validation process (trainers, then courses) that guarantees
the quality of published content. The architecture, decoupled into four
components --- a Spring Boot backend, a React Native mobile application, a
Next.js dashboard, and a FastAPI recommendation microservice --- made it
possible to handle each business concern independently while keeping a
single database as the source of truth.

The Scrum approach, split into four releases and eight sprints, made it
possible to deliver tested functional increments at every iteration, from
the authentication foundation to fine-grained financial oversight
(centralized pricing, transparent daily trainer compensation). Beyond the
core CRUD and workflow features, the platform integrates a genuinely
non-trivial machine learning component --- a hybrid content-based /
collaborative recommendation engine served through its own microservice,
with FAISS-accelerated similarity search, zero-downtime model refresh, and
a bounded per-user result cache --- as well as an AI-generated,
multi-language social feed, a role-based security model with
enumeration-resistant authentication endpoints and server-verified payment
webhooks, and a fully internationalized, RTL-aware mobile interface. Each
of these concerns was tested end to end, either through automated
black-box scripts (recommendation engine) or scenario-driven manual
validation against the product backlog's acceptance criteria (mobile
application and admin dashboard), before being accepted by LeanConsulting
at the close of its corresponding sprint.

Several avenues for future development emerge: social login (Google,
Apple, LinkedIn), real-time messaging via WebSocket, production release on
mobile stores through a containerized deployment and an EAS build with
native push notifications, a continuous-integration pipeline running the
existing test suites automatically on every change, and enriching the
recommendation engine with more advanced collaborative filtering
techniques such as matrix factorization.

% ============================================================
\begin{thebibliography}{9}
\addcontentsline{toc}{chapter}{Bibliography}

\bibitem{scrum} K.~Schwaber, J.~Sutherland, \emph{The Scrum Guide}, 2020.
\bibitem{uml} G.~Booch, J.~Rumbaugh, I.~Jacobson, \emph{The Unified Modeling
Language User Guide}, Addison-Wesley, 2005.
\bibitem{spring} Spring Team, \emph{Spring Boot Reference Documentation},
\url{https://docs.spring.io/spring-boot/}.
\bibitem{reactnative} Expo Team, \emph{Expo Documentation},
\url{https://docs.expo.dev/}.
\bibitem{nextjs} Vercel, \emph{Next.js Documentation},
\url{https://nextjs.org/docs}.
\bibitem{konnect} Konnect, \emph{Konnect API Documentation},
\url{https://developers.konnect.network/}.
\bibitem{sklearn} F.~Pedregosa et al., \emph{Scikit-learn: Machine Learning
in Python}, JMLR 12, 2011, \url{https://scikit-learn.org/}.
\bibitem{faiss} J.~Johnson, M.~Douze, H.~J\'egou, \emph{Billion-scale
similarity search with GPUs}, IEEE Transactions on Big Data, 2019,
\url{https://github.com/facebookresearch/faiss}.
\bibitem{gemini} Google, \emph{Gemini API Documentation},
\url{https://ai.google.dev/gemini-api/docs}.
\bibitem{expopush} Expo Team, \emph{Expo Push Notifications},
\url{https://docs.expo.dev/push-notifications/overview/}.
\bibitem{owasp} OWASP Foundation, \emph{OWASP Top 10 --- Web Application
Security Risks}, \url{https://owasp.org/www-project-top-ten/}.
\bibitem{i18next} i18next, \emph{i18next Documentation},
\url{https://www.i18next.com/}.

\end{thebibliography}

\end{document}